\documentclass[12pt]{article}
\usepackage[a4paper,margin=1in]{geometry}
\usepackage[T1]{fontenc}
\usepackage{lmodern,microtype}
\usepackage{amsmath,amssymb,amsthm,mathtools}
\usepackage{booktabs,array,enumitem,xcolor}
\usepackage[numbers,sort&compress]{natbib}
\usepackage[colorlinks=true,linkcolor=blue!55!black,citecolor=blue!55!black,urlcolor=blue!55!black]{hyperref}
\linespread{1.07}

\newtheorem{theorem}{Theorem}[section]
\newtheorem{proposition}[theorem]{Proposition}
\newtheorem{corollary}[theorem]{Corollary}
\theoremstyle{remark}
\newtheorem{remark}[theorem]{Remark}
\newcommand{\cH}{\mathcal H}
\newcommand{\cK}{\mathcal K}
\newcommand{\cL}{\mathcal L}
\newcommand{\rank}{\operatorname{rank}}
\newcommand{\vecop}{\operatorname{vec}}
\newcommand{\norm}[1]{\left\lVert#1\right\rVert}

\title{The Frobenius Left-Multiplication Obstruction\\
\large Why a Physical Temporal Packet Is Not a Full Ambient Riesz Shell}
\author{Liang Wang}
\date{July 2026}

\begin{document}
\maketitle

\begin{abstract}
RH-189--RH-191 formulate an oblique Feshbach problem on the matrix-valued
physical orbit.  Before computing the complement inverse, one must identify
the ambient operator correctly.  If $A\in\mathbb C^{n\times n}$ and the
state is $X\in\mathbb C^{n\times m}$, the dynamics is
$\cL_A X=AX$.  Under column-major vectorization,
\[
 \cL_A\simeq I_m\otimes A.
\]
Hence
\[
 \det(zI-\cL_A)=\det(zI-A)^m,
 \qquad
 \rank P_\Gamma(\cL_A)=m\,\rank P_\Gamma(A).
\]
Every base eigenvalue is repeated once per source column.  Therefore the
literal full-Frobenius Riesz target used informally in the preceding packet
discussion is type-wrong: a rank-one root shell or rank-four packet shell
cannot equal a complete ambient Riesz projection when $m=64,128,$ or $256$.

Combined with the exact RH-189 count
$N_{\cL}=N_K+N_D$, the multiplicity law forces
$N_D=m-1$ around a packet root whenever that root matches one simple base
eigenvalue.  If it matches no base eigenvalue, the uniform Schur homotopy
must fail.  A 126-window, 416-root audit finds no complement-free compatible
case; the forced one-root complement burden ranges from 63 to 255.

This is a structural correction, not a rejection of the numerical root
signal.  The source-cyclic space generated by the single matrix $S$ has
dimension at most $n$ and need not carry the factor $m$.  It is the correct
next setting for a low-rank source-relative spectral channel.  Gate A and
all Hilbert--P\'olya downstream claims remain open.
\end{abstract}

\section{The type hidden by the notation}

The RH-185 orbit is written $S,AS,A^2S,\ldots$, with
$S\in\mathbb C^{n\times m}$.  Treating each matrix as one Frobenius vector is
legitimate, but it does not turn left multiplication into a generic
$nm$-dimensional operator.  It preserves the source-column coordinate
exactly.  This elementary observation changes the Riesz-count target before
any numerical inversion is attempted.

Equip
\begin{equation}\label{eq:frobenius-space}
 \cH_F=\mathbb C^{n\times m}
\end{equation}
with the Frobenius inner product
$\langle X,Y\rangle_F=\operatorname{tr}(X^*Y)$, and define
\begin{equation}\label{eq:left-multiplication}
 \cL_A X=AX.
\end{equation}
The issue is not subtle pseudospectral behavior.  It is an exact direct-sum
decomposition, using standard finite matrix and Riesz-calculus identities
\citep{HornJohnson1991,Kato1995}.

\section{Kronecker equivalence}

Let $E_j:\mathbb C^n\to\cH_F$ place a vector in column $j$ and put zeros in
the other columns.  Then
\begin{equation}\label{eq:reducing-columns}
 \cH_F=\bigoplus_{j=1}^m E_j\mathbb C^n,
 \qquad
 \cL_A E_j=E_j A.
\end{equation}
The sum is orthogonal in the Frobenius norm.

\begin{theorem}[Unitary direct-sum form]\label{thm:direct-sum}
Column-major vectorization is unitary from $\cH_F$ to
$\mathbb C^m\otimes\mathbb C^n$ and satisfies
\begin{equation}\label{eq:kron}
 \vecop(AX)=(I_m\otimes A)\vecop(X).
\end{equation}
Consequently $\cL_A$ is unitarily equivalent to the orthogonal direct sum of
$m$ copies of $A$.
\end{theorem}

\begin{proof}
Write $X=[x_1,\ldots,x_m]$.  Then
$AX=[Ax_1,\ldots,Ax_m]$, while column-major vectorization stacks precisely
these columns.  The Frobenius norm is
$\norm X_F^2=\sum_j\norm{x_j}_2^2$, so the identification is unitary.
\end{proof}

\begin{corollary}[Characteristic and minimal polynomials]\label{cor:polynomials}
For every $z\in\mathbb C$,
\begin{equation}\label{eq:char-power}
 \det(zI_{nm}-\cL_A)=\det(zI_n-A)^m.
\end{equation}
The minimal polynomial of $\cL_A$ equals the minimal polynomial of $A$.
Thus algebraic multiplicities are multiplied by $m$, although Jordan-block
sizes are not enlarged.
\end{corollary}

\begin{proof}
The determinant of a block diagonal matrix is the product of its block
determinants.  A polynomial annihilates the direct sum if and only if it
annihilates every copy of $A$.
\end{proof}

The distinction in the last sentence is useful.  The obstruction is not an
artificial lengthening of Jordan chains; it is an exact replication of every
chain across source columns.

\section{Resolvents and Riesz projectors}

For $z\notin\sigma(A)$, direct calculation gives
\begin{equation}\label{eq:resolvent-action}
 (zI-\cL_A)^{-1}X=(zI-A)^{-1}X.
\end{equation}

\begin{proposition}[No resolvent-norm penalty]\label{prop:resolvent}
In the induced Frobenius norm,
\begin{equation}\label{eq:resolvent-norm}
 \norm{(zI-\cL_A)^{-1}}_{F\to F}
 =\norm{(zI-A)^{-1}}_{2\to2}.
\end{equation}
\end{proposition}

\begin{proof}
The direct-sum norm is the maximum of the block norms, giving the upper
bound.  Equality follows by placing a maximizing right singular vector of
$(zI-A)^{-1}$ in one column and zeros elsewhere.
\end{proof}

Thus the Frobenius lift does not worsen the resolvent norm.  It changes the
spectral count.

Let $\Gamma$ be a positively oriented contour in the resolvent set of $A$,
and let
\begin{equation}\label{eq:base-riesz}
 P_\Gamma(A)=\frac{1}{2\pi i}\int_\Gamma(zI-A)^{-1}\,dz.
\end{equation}

\begin{theorem}[Frobenius Riesz multiplicity]\label{thm:riesz-rank}
The ambient Riesz projector acts columnwise,
\begin{equation}\label{eq:projector-action}
 P_\Gamma(\cL_A)X=P_\Gamma(A)X,
\end{equation}
and
\begin{equation}\label{eq:rank-law}
 \boxed{\rank P_\Gamma(\cL_A)
 =m\,\rank P_\Gamma(A).}
\end{equation}
Equivalently, the algebraic eigenvalue count obeys
$N_{\cL_A}(\Gamma)=mN_A(\Gamma)$.
\end{theorem}

\begin{proof}
Insert \eqref{eq:resolvent-action} into the contour integral.  The range is
the direct sum of $m$ copies of $\operatorname{Ran}P_\Gamma(A)$, so dimensions
multiply.
\end{proof}

\section{Consequence for the RH-189 Feshbach count}

RH-189 proves that a uniform Schur margin for packet/complement coordinates
implies
\begin{equation}\label{eq:feshbach-count}
 N_{\cL_A}(\Gamma)=N_K(\Gamma)+N_D(\Gamma).
\end{equation}
Combining this with Theorem~\ref{thm:riesz-rank} yields an arithmetic law
that must be checked before estimating norms.

\begin{corollary}[Complement congruence]\label{cor:congruence}
Under the hypotheses of the RH-189 count theorem,
\begin{equation}\label{eq:congruence}
 N_D(\Gamma)=mN_A(\Gamma)-N_K(\Gamma),
 \qquad
 N_D(\Gamma)\equiv-N_K(\Gamma)\pmod m.
\end{equation}
In particular, a complement-free count is possible only if $m$ divides
$N_K(\Gamma)$.
\end{corollary}

For a contour containing one simple packet eigenvalue, $N_K=1$.  There are
only two cases:
\begin{enumerate}[leftmargin=2em]
\item If $N_A=0$, equation \eqref{eq:congruence} asks for $N_D=-1$, which is
 impossible.  Hence at least one assumption of the Schur homotopy fails.
\item If $N_A=1$, the complement count is $N_D=m-1$.  The packet may locate
 the base eigenvalue, but it cannot be the complete ambient Riesz shell.
\end{enumerate}

The same issue remains for the union of four disjoint root contours.  If
each contains one simple base eigenvalue, the full ambient count is $4m$,
not four, and the complement carries $4(m-1)$ dimensions.

\section{Physical arithmetic audit}

The RH-185 data contain 126 temporal windows and 416 packet roots.  The
matrix-state widths are:
\begin{center}
\begin{tabular}{crrr}
\toprule
$\sigma$ & representative $m$ & one-root full count & forced complement\\
\midrule
$0.04$ & $64$ & $64$ & $63$\\
$0.02$ & $128$ & $128$ & $127$\\
$0.01$ & $256$ & $256$ & $255$\\
\bottomrule
\end{tabular}
\end{center}
The left and right channels can have different base dimensions $n$, but
they use the same width at a fixed scale.  The multiplicity law depends on
$m$, not on $n$.

No one-root count and no complete length-three or length-four packet count
is divisible by its physical width.  Thus all 126 literal complement-free
targets are arithmetically impossible.  This conclusion is exact once the
state type is fixed; it does not depend on floating eigenvalue accuracy,
contour meshes, or an inverse norm estimate.

\section{What survives: the single-seed cyclic space}

The direct-sum theorem does not say that every physically relevant
subspace has dimension multiplied by $m$.  Let
\begin{equation}\label{eq:cyclic-space}
 \cK_S=\operatorname{span}\{S,AS,A^2S,\ldots\}\subset\cH_F.
\end{equation}
This is generated by one vector of the Frobenius Hilbert space.  It is
$\cL_A$-invariant and has dimension at most the degree of the minimal
polynomial of $A$, hence at most $n$.  Its restriction is cyclic and can have
simple eigenvalue counts.

The RH-185 right temporal synthesis lies in $\cK_S$ by construction.  Thus
the numerical packet may still be identifying one source-excited direction
from each $m$-fold ambient eigenspace.  The error was to identify that
source-relative direction with the complete ambient eigenspace.

An observation seed supplies a further quotient: modes invisible to the
input-output moments can be removed.  These two reductions are standard in
realization theory, but they must be implemented without silently restoring
the full-Frobenius count.

\section{Revised decision tree}

The next calculation should not invert a complement of dimension $nm-L$
and then hope for a rank-$L$ ambient shell.  The corrected sequence is:
\begin{enumerate}[leftmargin=2em]
\item construct the invariant source-cyclic restriction of $\cL_A$;
\item prove that the temporal right packet is contained in it and restrict
 the left functionals to it;
\item identify the source-observable spectral modes selected by the
 physical observation;
\item only then define packet/complement coordinates and Riesz counts in
 that reduced state type.
\end{enumerate}

If the temporal roots fail to match the spectrum of the cyclic restriction,
the local candidate is rejected.  If they match, one obtains a genuine
source-channel spectral packet, but not yet a canonical all-level cloud.

\section{Boundary of the result}

This paper proves a no-go theorem for one literal interpretation only:
\begin{equation*}
 \text{rank-one/rank-four temporal packet}
 =\text{complete Riesz shell of }\cL_A\text{ on }\mathbb C^{n\times m}.
\end{equation*}
It does not reject the RH-185 eigenvalue approximation, a source-cyclic
restriction, a controllable/observable minimal realization, or a different
physical operator that genuinely mixes source columns.

The distinction is important for the larger program.  Gate A requires an
intrinsic and transportable physical spectral interface.  Correcting the
ambient multiplicity is necessary for that interface, but it does not
construct it.  Gates B--E, a self-adjoint Hilbert--P\'olya operator, zeta-zero
identification, and the Riemann hypothesis are untouched.

\section{Reproducibility}

The accompanying implementation tests the Kronecker action, determinant
power, resolvent norm, and count arithmetic on independent matrices.  The
physical audit reads the frozen RH-185 records and stores every width,
ambient dimension, packet length, and forced complement burden.  It contains
no fitted threshold and no spectral matching step.

The full command sequence is recorded in the README.  A publication
manifest hashes source, tests, frozen inputs referenced by name, result JSON,
and the compiled paper.

\bibliographystyle{plainnat}
\bibliography{references}
\end{document}
